\section*{第 \ref{ch:g11:fossil-fuel-geology} 章——煤、石油和天然气的地质学}

\begin{solution}{exo:g11:fossil-fuel-geology:1}
死亡的浮游生物沉降在贫氧的海底；有机质
在泥中转化为\omterm{def:g11:fossil-fuel-geology:source-rock}{干酪根}；被更年轻的沉积物掩埋；在
\omterm{def:g11:fossil-fuel-geology:oil-window}{生油窗}内\omterm{def:g11:fossil-fuel-geology:oil-window}{成熟}；发生\omterm{def:g11:fossil-fuel-geology:migration}{运移}；石油聚集在\omterm{def:g11:fossil-fuel-geology:reservoir}{盖层}之下。
\end{solution}

\begin{solution}{exo:g11:fossil-fuel-geology:2}
(1) \omterm{def:g11:fossil-fuel-geology:source-rock}{烃源岩}。(2) \omterm{def:g11:fossil-fuel-geology:reservoir}{储集岩}。(3) \omterm{def:g11:fossil-fuel-geology:reservoir}{盖层}。(4) \omterm{def:g11:fossil-fuel-geology:reservoir}{盖层}。(5) \omterm{def:g11:fossil-fuel-geology:reservoir}{储集岩}。
\end{solution}

\begin{solution}{exo:g11:fossil-fuel-geology:3}
顶部：$(60 - 10)/25 = \qty{2.0}{km}$；底部：$(120 - 10)/25 = \qty{4.4}{km}$。
\omterm{def:g11:fossil-fuel-geology:oil-window}{生油窗}的厚度为 $60/25 = \qty{2.4}{km}$。
\end{solution}

\begin{solution}{exo:g11:fossil-fuel-geology:4}
(1) 泥炭、褐煤、烟煤、无烟煤。(2) 无烟煤，
碳含量为 \qty{86}{\percent} 至 \qty{97}{\percent}。(3) 无烟煤，埋深超过约 \qty{5}{km} 且
温度超过 \qty{150}{\degreeCelsius}。
\end{solution}

\begin{solution}{exo:g11:fossil-fuel-geology:5}
(1) 流体按密度自行分层：天然气最轻，其次为石油，
最后为水。(2) 石油和天然气的密度小于水，因而向上\omterm{def:g11:fossil-fuel-geology:migration}{运移}；只有
它们无法穿过的岩层才能阻挡其\omterm{def:g11:fossil-fuel-geology:migration}{运移}。泥岩（或页岩）和盐岩。
\end{solution}

\begin{solution}{exo:g11:fossil-fuel-geology:6}
(1) A：$(60 - 15)/20 = \qty{2.25}{km}$，$(120 - 15)/20 = \qty{5.25}{km}$。
B：$45/40 \approx \qty{1.13}{km}$，$105/40 \approx \qty{2.63}{km}$。
(2) 在 \qty{2}{km} 处：在 A 中，$15 + 40 = \qty{55}{\degreeCelsius}$，尚未
\omterm{def:g11:fossil-fuel-geology:oil-window}{成熟}，无石油；在 B 中，$15 + 80 = \qty{95}{\degreeCelsius}$，位于
\omterm{def:g11:fossil-fuel-geology:oil-window}{生油窗}内，可生成石油。(3) 在盆地 A 中，石油仅在
\qty{2.25}{km} 以下才会生成。
\end{solution}

\begin{solution}{exo:g11:fossil-fuel-geology:7}
(1) $11 \times 3 = \qty{33}{m}$ 泥炭。(2) $\qty{33000}{mm}$ 以
每年 \qty{1}{mm} 计，约需 \num{33000} 年。(3) 泥炭只有在
停滞水体之下、与氧气隔绝时才能保存；随着泥炭不断堆积，
地面必须以相同速度下沉，才能使其表面保持被水浸没的状态。
\end{solution}

\begin{solution}{exo:g11:fossil-fuel-geology:8}
(1) $20 + 30 \times 1.2 = \qty{56}{\degreeCelsius}$；
$20 + 30 \times 2.3 = \qty{89}{\degreeCelsius}$；
$20 + 30 \times 4.0 = \qty{140}{\degreeCelsius}$。
(2) 1.2 亿年前，尚未\omterm{def:g11:fossil-fuel-geology:oil-window}{成熟}；9500 万年前，位于
\omterm{def:g11:fossil-fuel-geology:oil-window}{生油窗}内；2000 万年前，位于\omterm{def:g11:fossil-fuel-geology:oil-window}{生气窗}内。
\end{solution}

\begin{solution}{exo:g11:fossil-fuel-geology:9}
(1) $2580 - 2400 = \qty{180}{m}$。(2) 天然气充满
$2450 - 2400 = \qty{50}{m}$，石油充满 $2580 - 2450 = \qty{130}{m}$。(3) 它们在
溢出点处溢出，该点位于最深闭合等高线之下，随后继续向上
\omterm{def:g11:fossil-fuel-geology:migration}{运移}，进入另一\omterm{def:g11:fossil-fuel-geology:trap}{圈闭}或到达地表。
\end{solution}

\begin{solution}{exo:g11:fossil-fuel-geology:10}
(1) \omterm{def:g11:fossil-fuel-geology:oil-window}{生油窗}顶部位于 \qty{3}{km} 处：$(60 - 20)/3 \approx
\qty{13}{\degreeCelsius/km}$；底部：$(120 - 20)/3 \approx
\qty{33}{\degreeCelsius/km}$。(2) 在 \qty{3}{km} 处温度为
$20 + 75 = \qty{95}{\degreeCelsius}$ 至 $20 + 102 = \qty{122}{\degreeCelsius}$：
位于\omterm{def:g11:fossil-fuel-geology:oil-window}{生油窗}内，对于最陡的\omterm{def:g11:lithosphere-asthenosphere:geotherm}{地温梯度}则位于其底部。该规则
是合理的：在大多数盆地中，\omterm{def:g11:fossil-fuel-geology:source-rock}{烃源岩}上覆 \qty{3}{km} 厚的岩层时，
已开始生成石油。
\end{solution}

\begin{solution}{exo:g11:fossil-fuel-geology:11}
(1) \qty{30}{\percent}：褐煤；\qty{70}{\percent}：烟煤；
\qty{92}{\percent}：无烟煤。(2) 无烟煤，埋深超过约
\qty{5}{km}、温度超过 \qty{150}{\degreeCelsius}。(3) 黑色，玻璃质。
\end{solution}

\begin{solution}{exo:g11:fossil-fuel-geology:12}
(1) $1100/25 = \qty{44}{m}$ 每百万年。(2) $0.044 \times 30
\approx \qty{1.3}{\degreeCelsius}$ 每百万年。(3) \omterm{def:g11:fossil-fuel-geology:oil-window}{生油窗}跨越
\qty{60}{\degreeCelsius}：$60/1.32 \approx 45$ 百万年。(4) 埋藏
速率在 9500 万年前之后减慢（每百万年约 24 米，随后为 23 米和
\qty{20}{m}）：\omterm{def:g11:fossil-fuel-geology:source-rock}{烃源岩}升温变慢。
\end{solution}

\begin{solution}{exo:g11:fossil-fuel-geology:13}
(1) $(80 - 51)/(117 - 51) = 29/66 \approx 0.44$：约 \qty{44}{\percent}。
(2) 它在其上方形成任何\omterm{def:g11:fossil-fuel-geology:trap}{圈闭}之前已向上\omterm{def:g11:fossil-fuel-geology:migration}{运移}：逸散至
地表并被氧化，或聚集在其他\omterm{def:g11:fossil-fuel-geology:trap}{圈闭}中。(3) 4000 万年前
\omterm{def:g11:fossil-fuel-geology:source-rock}{烃源岩}已离开\omterm{def:g11:fossil-fuel-geology:oil-window}{生油窗}：只生成天然气，
因此背斜只能容纳来自该\omterm{def:g11:fossil-fuel-geology:source-rock}{烃源岩}的天然气。
\end{solution}

\begin{solution}{exo:g11:fossil-fuel-geology:14}
(1) 顶部：$(60 - 20)/40 = \qty{1.0}{km}$；底部：$(120 - 20)/40 =
\qty{2.5}{km}$。(2) \qty{1.0}{km} 介于 \qty{0.4}{km}（1.5 亿年前
）与 \qty{1.2}{km}（1.2 亿年前）之间：$150 - 30 \times 0.6/0.8 = 127.5$
百万年前。\qty{2.5}{km} 介于 \qty{2.3}{km}（9500 万年前）与
\qty{3.0}{km}（6600 万年前）之间：$95 - 29 \times 0.2/0.7 \approx 86.7$
百万年前。(3) 无石油：\omterm{def:g11:fossil-fuel-geology:source-rock}{烃源岩}在背斜形成之前已离开\omterm{def:g11:fossil-fuel-geology:oil-window}{生油窗}，
其生成的石油已\omterm{def:g11:fossil-fuel-geology:migration}{运移}走。它只能容纳 8000 万年前之后
生成的天然气。
\end{solution}

\begin{solution}{exo:g11:fossil-fuel-geology:15}
(1) 至少 \qty{5}{km}。(2) $15 + 30 \times 5 = \qty{165}{\degreeCelsius}$，
高于 \qty{150}{\degreeCelsius}：结果一致。(3) 至少 \qty{5}{km} 厚的
岩层已被抬升并遭受剥蚀。$\qty{5e6}{mm}$ 在至多
2.99 亿年内完成，平均每年至少为 $5 \times 10^6/2.99 \times 10^8
\approx \qty{0.017}{mm}$。
\end{solution}

\begin{solution}{pb:g11:fossil-fuel-geology:1}
\textbf{1.} \omterm{def:g11:fossil-fuel-geology:source-rock}{烃源岩}：黑色页岩；\omterm{def:g11:fossil-fuel-geology:reservoir}{储集岩}：多孔砂岩；
\omterm{def:g11:fossil-fuel-geology:reservoir}{盖层}：其上方的泥岩。

\textbf{2.} 沉积于缺氧海底的浮游生物
得以避免腐烂，并被埋入泥中。

\textbf{3.} \omterm{def:g11:fossil-fuel-geology:source-rock}{烃源岩}形成于侏罗纪（165 位于 201.4 与
143.1 之间）。砂岩始于侏罗纪，主要沉积于
早白垩世。

\textbf{4.} 分散于沉积岩中的有机物，不溶于
有机溶剂。只有加热至约
\qty{60}{\degreeCelsius}时才会生成石油，这需要埋藏作用。

\textbf{5.} $400 + 800 + 1100 + 700 + 700 + 300 - 500 = \qty{3500}{m}$.

\textbf{6.} \qty{0.4}{km}, \qty{1.2}{km}, \qty{2.3}{km}, \qty{3.0}{km},
\qty{3.7}{km}、\qty{4.0}{km}；现今为 \qty{3.5}{km}。

\textbf{7.} $400/15 \approx 26.7$; $800/30 \approx 26.7$; $1100/25 = 44$;
$700/29 \approx 24.1$；$700/31 \approx 22.6$；$300/15 = 20$（每
百万年 \unit{m}）。

\textbf{8.} 2000 万年前，埋深为 \qty{4.0}{km}。

\textbf{9.} 盆地被抬升，其顶部被剥蚀掉
\qty{500}{m} 厚的岩层。

\textbf{10.} 砂岩底部沉积于 1.5 亿年前，
位于 \qty{400}{m} 厚的页岩与泥岩之上：$3.5 - 0.4 = \qty{3.1}{km}$。

\textbf{11.} $T = 20 + 30 z$: \qty{32}{\degreeCelsius}, \num{56}, \num{89},
\num{110}、\num{131}、\qty{140}{\degreeCelsius}；现今为
\qty{125}{\degreeCelsius}。

\textbf{12.} $40/30 \approx \qty{1.33}{km}$ 和 $100/30 \approx
\qty{3.33}{km}$.

\textbf{13.} 介于 1.2 亿年前（\qty{56}{\degreeCelsius}）与
9500 万年前（\qty{89}{\degreeCelsius}）之间。

\textbf{14.} 温度在 2500 万年内升高了 33 度，所用时间
约为 $25 \times 4/33 \approx 3.0$ 百万年，升高 4 度，
即自 \num{56} 升至 \qty{60}{\degreeCelsius}，约在 $120 - 3.0 = 117$ 百万
年前。

\textbf{15.} 介于 6600 万年前（\qty{110}{\degreeCelsius}）与
3500 万年前（\qty{131}{\degreeCelsius}）之间：$66 - 31 \times 10/21 \approx 51.2$ 百万
年前。

\textbf{16.} 2000 万年前为 \qty{140}{\degreeCelsius}。
现今为 \qty{125}{\degreeCelsius}，已处于\omterm{def:g11:fossil-fuel-geology:oil-window}{生气窗}内：
主要生成甲烷。

\textbf{17.} 其上方无\omterm{def:g11:fossil-fuel-geology:trap}{圈闭}，因而向上迁移并逸出，
至地表或其他\omterm{def:g11:fossil-fuel-geology:trap}{圈闭}中。

\textbf{18.} $(80 - 51.2)/(117.0 - 51.2) \approx 0.44$.

\textbf{19.} 顶部为气，其次为油，最下为水。

\textbf{20.} \omterm{def:g11:fossil-fuel-geology:source-rock}{烃源岩}约于 \textbf{1.17
亿年前}进入\omterm{def:g11:fossil-fuel-geology:oil-window}{生油窗}。
\end{solution}
